% from ega2-1-fr.tex
\oldpage[II]{158}
est une $S$-algèbre graduée, quand on prend pour degré de $fz^n$ ($n\geq0$) où
$f$ est homogène,
\begin{equation}
\label{II.8.2.2.2-fr}
  \deg(fz^n)=n+\deg f.
\tag{8.2.2.2}
\end{equation}
\end{env}

\begin{lemma}[8.2.3]
\label{II.8.2.3-fr}
\begin{enumerate}
\item[\rm(i)] On a des isomorphismes canoniques d'anneaux (non gradués)
\begin{equation}
\label{II.8.2.3.1-fr}
  \widehat{S}_{(z)}\xrightarrow{\sim}
  \widehat{S}/(z-1)\widehat{S}\xrightarrow{\sim}S.
\tag{8.2.3.1}
\end{equation}
\item[\rm(ii)] On a un isomorphisme canonique d'anneaux (non gradués)
\begin{equation}
\label{II.8.2.3.2-fr}
  \widehat{S}_{(f)}\xrightarrow{\sim}S_f^{\leq}
\tag{8.2.3.2}
\end{equation}
pour tout $f\in S_d$, avec $d>0$.
\end{enumerate}
\end{lemma}

Le premier des isomorphismes de (\hyperref[II.8.2.3.1-fr]{8.2.3.1}) a été
défini dans (\hyperref[II.2.2.5-fr]{2.2.5}) et le second est trivial~;
l'isomorphisme $\widehat{S}_{(z)}\xrightarrow{\sim}S$ ainsi défini fait donc
correspondre à $xz^n/z^{n+k}$, où $\deg(x)=k$ ($k\geq-n$) l'élément $x$.
L'homomorphisme (\hyperref[II.8.2.3.2-fr]{8.2.3.2}) fait correspondre à
$xz^n/f^k$, où $\deg(x)=kd-n$, l'élément $x/f^k$, de degré $-n$ dans
$S_f^{\leq}$, et il est encore évident qu'on a bien ici un isomorphisme.

\begin{env}[8.2.4]
\label{II.8.2.4-fr}
Soit $M$ un $S$-module gradué. Il est clair que le $S$-module
\begin{equation}
\label{II.8.2.4.1-fr}
  \widehat{M}=M\otimes_S\widehat{S}=M\otimes_S S[z]
\tag{8.2.4.1}
\end{equation}
est somme directe des $S$-modules $M\otimes_S Sz^n$, donc des groupes abéliens
$M_k\otimes_S Sz^n$ ($k\in\mathbf{Z}$, $n\geq0$)~; on définit sur
$\widehat{M}$ une structure de $\widehat{S}$-module gradué en prenant
\begin{equation}
\label{II.8.2.4.2-fr}
  \deg(x\otimes z^n)=n+\deg x
\tag{8.2.4.2}
\end{equation}
pour tout $x$ homogène dans $M$. Nous laissons au lecteur le soin de démontrer
l'analogue de (\hyperref[II.8.2.3-fr]{8.2.3})~:
\end{env}

\begin{lemma}[8.2.5]
\label{II.8.2.5-fr}
\begin{enumerate}
\item[\rm(i)] On a un di-isomorphisme canonique de modules (non gradués)
\begin{equation}
\label{II.8.2.5.1-fr}
  \widehat{M}_{(z)}\xrightarrow{\sim}M.
\tag{8.2.5.1}
\end{equation}
\item[\rm(ii)] Pour tout $f\in S_d$ ($d>0$), on a un di-isomorphisme de
modules (non gradués)
\begin{equation}
\label{II.8.2.5.2-fr}
  \widehat{M}_{(f)}\xrightarrow{\sim}M_f^{\leq}.
\tag{8.2.5.2}
\end{equation}
\end{enumerate}
\end{lemma}

\begin{env}[8.2.6]
\label{II.8.2.6-fr}
Soit $S$ un anneau gradué à degrés \emph{positifs}, et considérons dans $S$
la suite décroissante d'idéaux gradués
\begin{equation}
\label{II.8.2.6.1-fr}
  S_{[n]}=\bigoplus_{m\geq n}S_m\qquad(n\geq0)
\tag{8.2.6.1}
\end{equation}
(en particulier on a $S_{[0]}=S$, $S_{[1]}=S_+$). Comme il est clair que
$S_{[m]}S_{[n]}\subset S_{[m+n]}$, on peut définir un \emph{anneau gradué}
$S^\natural$ en prenant
\begin{equation}
\label{II.8.2.6.2-fr}
  S^\natural=\bigoplus_{n\geq0}S_n^\natural
  \quad\text{avec}\quad S_n^\natural=S_{[n]}.
\tag{8.2.6.2}
\end{equation}
$S_0^\natural$ est donc l'anneau $S$ considéré comme anneau \emph{non gradué},
et $S^\natural$ est par suite une $S_0^\natural$-algèbre. Pour tout élément
homogène $f\in S_d$ ($d>0$), nous désignerons par $f^\natural$ l'élément
$f$ considéré comme appartenant à $S_{[d]}=S_d^\natural$. Avec ces
notations~:
\end{env}

